\chapter{Init System}
\label{ch:initsys}

The init system is the first userspace process (PID 1) started by the
kernel. It is responsible for bringing up the system, managing services,
and handling shutdown. Veilbox uses BusyBox init — a minimal, predictable
init implementation.

\section{PID 1 Responsibilities}

Every init system must handle:

\begin{itemize}[nosep]
  \item \textbf{Orphaned processes}: When a child's parent dies, the child
        is reparented to PID 1. Init must \texttt{wait()} for orphaned
        children to prevent zombie processes.
  - \textbf{Signals}: PID 1 receives special signal handling. The kernel
        does not deliver SIGKILL to PID 1 (it would be ignored).
  - \textbf{Shutdown sequence}: On reboot or poweroff, init must unmount
        filesystems, kill processes, and signal the kernel to power off.
  - \textbf{Service management}: Init spawns system services and restarts
        them if they die (crashed services get respawned).
\end{itemize}

\section{BusyBox Init Operation}

BusyBox init follows a simple, deterministic loop:

\begin{enumerate}[nosep]
  \item \textbf{Parse /etc/inittab}: Read configuration file into an
        internal list of actions.
  \item \textbf{Execute sysinit}: Run the single command specified with
        the \texttt{sysinit} action. Wait for completion.
  \item \textbf{Spawn respawn processes}: Fork and exec each process
        marked with the \texttt{respawn} or \texttt{askfirst} action.
        Store PIDs in a tracking table.
  \item \textbf{Wait loop}: Call \texttt{waitpid()} in a loop. When a
        child dies, check if it was a \texttt{respawn} process. If so,
        re-fork and exec it. If it was a \texttt{once} process, do
        nothing (it ran and exited).
  \item \textbf{Signal handling}: On SIGTERM, SIGQUIT, SIGUSR1, or
        SIGUSR2, execute the corresponding \texttt{shutdown} or
        \texttt{restart} action from inittab.
\end{enumerate}

\section{The /etc/inittab File}

BusyBox inittab format:

\begin{lstlisting}
<id>:<runlevels>:<action>:<process>
\end{lstlisting}

\begin{longtable}{|p{2.5cm}|p{4cm}|p{6cm}|}
\hline
\textbf{Field} & \textbf{Description} & \textbf{Veilbox Usage} \\
\hline
\texttt{<id>} & TTY device (optional). & \texttt{tty1}, \texttt{ttyS0} \\
\texttt{<runlevels>} & Ignored by BusyBox. & Empty \\
\texttt{<action>} & One of: sysinit, respawn, askfirst, once, wait, restart, shutdown, ctrlaltdel. & \texttt{sysinit}, \texttt{respawn}, \texttt{shutdown} \\
\texttt{<process>} & Full path to executable. & \texttt{/etc/init.d/rcS}, \texttt{/sbin/getty} \\
\hline
\end{longtable}

Veilbox inittab:

\begin{lstlisting}
::sysinit:/etc/init.d/rcS
tty1::respawn:/sbin/autologin 38400 tty1
ttyS0::respawn:/sbin/autologin 38400 ttyS0
::restart:/sbin/init
::shutdown:/bin/umount -a -r
::shutdown:/sbin/swapoff -a
\end{lstlisting}

\section{Boot Script: /etc/init.d/rcS}

The rcS script is the heart of system initialization. It runs
sequentially and blocks init until complete:

\begin{lstlisting}[language=bash]
#!/bin/sh
# Mount virtual filesystems
mount -t proc none /proc
mount -t sysfs none /sys
mount -t devtmpfs none /dev
mkdir -p /dev/pts
mount -t devpts none /dev/pts
mount -t tmpfs none /tmp

# Set hostname
hostname -F /etc/hostname

# Mount state partition
STATE_DEV=""
if [ -b /dev/vdb ]; then
    STATE_DEV=/dev/vdb
elif [ -b /dev/sdb ]; then
    STATE_DEV=/dev/sdb
fi
if [ -n "$STATE_DEV" ]; then
    mount "$STATE_DEV" /mnt/state
fi

# Configure network
ifconfig lo 127.0.0.1 up
ifconfig eth0 0.0.0.0 up
udhcpc -i eth0 -b

# Start services
syslogd -O /var/log/messages -s 64
containerd -c /etc/containerd/config.toml &
dropbear -F -R -p 22 &
\end{lstlisting}

\section{Shutdown Sequence}

When \texttt{poweroff} or \texttt{reboot} is issued:

\begin{enumerate}[nosep]
  \item BusyBox init receives SIGTERM.
  \item init sends SIGTERM to all processes (except PID 1).
  \item init runs \texttt{shutdown} actions from inittab: unmount
        filesystems and deactivate swap.
  \item init sends SIGKILL to remaining processes.
  \item init syncs filesystems and calls \texttt{reboot(LINUX\_REBOOT\_CMD\_POWER\_OFF)}.
\end{enumerate}

\section{Autologin}

The \texttt{/sbin/autologin} wrapper checks if \texttt{veilbox.autologin}
is in the kernel command line:

\begin{lstlisting}[language=bash]
#!/bin/sh
if grep -q veilbox.autologin /proc/cmdline; then
    login -f root
else
    exec /sbin/getty "$@"
fi
\end{lstlisting}

\section{Comparing Init Systems}

\begin{longtable}{|p{3cm}|p{3.5cm}|p{3.5cm}|p{3.5cm}|}
\hline
\textbf{Feature} & \textbf{BusyBox init} & \textbf{systemd} & \textbf{OpenRC} \\
\hline
Binary size & 8 KB & 1.5 MB & 200 KB \\
Config files & /etc/inittab & .service units & /etc/init.d/ \\
Parallel startup & No & Yes & Partial \\
Dependency tracking & No & Yes & Yes \\
Socket activation & No & Yes & No \\
Timer support & No & Yes & No \\
Journal \& logging & No & journald & syslog-ng \\
\hline
\end{longtable}

\section{Adding New Services}

To add a new service to Veilbox:

\begin{enumerate}[nosep]
  \item Modify the rcS script to start the service:
\begin{lstlisting}[language=bash]
# /etc/init.d/rcS append
echo "Starting my-service..."
my-service --config /etc/my-service.conf &
\end{lstlisting}
  \item Add a configuration file:
\begin{lstlisting}[language=bash]
mkdir -p rootfs/etc
cat > rootfs/etc/my-service.conf << 'EOF'
# My service configuration
port = 8080
debug = false
EOF
\end{lstlisting}
  \item Download and install the binary:
\begin{lstlisting}[language=bash]
curl -Lo rootfs/usr/bin/my-service \
    https://example.com/my-service
chmod +x rootfs/usr/bin/my-service
\end{lstlisting}
  \item Regenerate initramfs and rebuild kernel:
\begin{lstlisting}[language=bash]
./build.sh --stage initramfs
./build.sh --stage kernel
\end{lstlisting}
  \item Test by booting and verifying the service:
\begin{lstlisting}
ps | grep my-service
my-service --version
\end{lstlisting}
\end{enumerate}

\section{System Rescue Mode}

If the system fails to boot normally, several rescue modes are available:

\subsection{Single-User Mode}

Boot with \texttt{single} on the kernel command line. The system boots
to a root shell without starting services:

\begin{lstlisting}
# In GRUB, edit the menu entry and add "single" to the linux line
linux /boot/vmlinuz console=ttyS0 console=tty1 single
\end{lstlisting}

\subsection{Direct Shell (No Init)}

Boot with \texttt{init=/bin/sh} to skip init entirely:

\begin{lstlisting}
linux /boot/vmlinuz console=ttyS0 init=/bin/sh
\end{lstlisting}

The root filesystem will be mounted but services will not start. The
shell has PID 1 and the root filesystem is mounted read-write.

\subsection{Read-Only Rootfs Recovery}

If the root filesystem is corrupted, boot with \texttt{ro} and manually
run fsck:

\begin{lstlisting}
# Not applicable to initramfs (tmpfs cannot corrupt)
# But for the state partition:
fsck.ext4 -f /dev/vdb
\end{lstlisting}

\section{Init Process Debugging}

BusyBox init can be run with verbose output for debugging:

\begin{lstlisting}
# Add to kernel command line
init=/sbin/init -v

# Or modify /etc/inittab to add debug output
::sysinit:/bin/sh -c "echo 'Starting init...'"
::sysinit:/etc/init.d/rcS
\end{lstlisting}
